\documentclass[10pt]{article}
\usepackage[margin=20mm]{geometry}
\usepackage{amsmath,amssymb,fontspec,unicode-math,xurl,hyperref}
\IfFontExistsTF{DejaVu Serif}{\setmainfont{DejaVu Serif}}{\setmainfont{Latin Modern Roman}}
\setmathfont{latinmodern-math.otf}
\setlength{\parindent}{0pt}\setlength{\parskip}{6pt}\sloppy
\begin{document}
Abstract: A missing tail term invalidates the proposed shuffle protocol’s unrestricted high-confidence sample-complexity argument, a Major concern tied directly to Theorem 1.4. An explicit failure event establishes the limitation for every fixed hidden constant. The formal quantile contract also permits valid datasets with no legal output, independently invalidating the as-written success guarantees (Major). Two Minor issues concern the released endpoint-selection code and the stated unary-encoding privacy budget. The intended adaptive LDP rate and adaptivity separation may survive a consistent output convention; those corrected results are not disproved here. Scientific decisions remain pending.\par
Date:\par
Title: Scientific Review of
Lightweight Protocols for Distributed
Private Quantile Estimation\par
\# Identity, material and assumptions\par
Paper 0114 / I25-43483 is the ICML 2025 paper by Anders Aamand, Fabrizio Boninsegna, Abigail Gentle, Jacob Imola and Rasmus Pagh, PMLR 267:27–58 [Retained latex: \textbackslash{}cite\{paper\}]. References below use printed PDF pages 1–32. The immutable supplied PDF and manifest were read and checked. The supplied source snapshot is pinned to `1301f6cdf662befb443dbfa7a94394b6a87bbdcf` [Retained latex: \textbackslash{}cite\{code\}]. It is selective: dataset and experimental-result binaries listed as excluded in the manifest were not reacquired.\par
The paper estimates one quantile of a fixed dataset \(X\in[B]^n\), with each user queried at most once. The adaptive local protocol privately reports one comparison bit; the shuffle protocol groups users into \(r=\log_2B\) disjoint batches. The review inspected the definitions, empirical-CDF drift proof, the noisy-search reduction, the potential argument, lower-bound reductions, shuffle analysis and relevant implementation paths. External sharp-capacity and nonadaptive lower-bound results were not fully rederived.\par
\# Core confidence dependence\par
Problem 1 (Major: the shuffle analysis drops a nonvanishing tail term).  Theorem 1.4, p. 3, and its restatement G.6, p. 25, give \[n=O\!\left((\alpha^{-2}+\varepsilon^{-2})\log B
 \sqrt{\log(1/\delta)\log(\log B/\beta)}\right)\] for success probability \(1-\beta\). The proof enters this bound through the high-\(\varepsilon_L\) simplification of Theorem G.3, p. 24. That simplification discards a term already present in Lemma G.2.\par
For a batch of \(s\) users, the displayed Bernstein calculation contains, up to constants and the two-sided failure allocation, \[\left[
 \frac{2e^{\varepsilon_L}}{\alpha^2(e^{\varepsilon_L}-1)^2}
 +\frac{2(e^{\varepsilon_L}+1)}{3\alpha(e^{\varepsilon_L}-1)}
 \right]\log(1/\beta').\] As \(\varepsilon_L\to\infty\), its second term tends to \(2\log(1/\beta')/(3\alpha)\); it is not \(O(e^{-\varepsilon_L}\alpha^{-2}\log(1/\beta'))\) uniformly. The proof of G.6 nevertheless starts with only the latter dependence, substitutes the amplification calibration, and obtains the square root in (2.1).\par
\#\#\#\# A complete protocol-level check.\par
Fix \[B=4,\quad r=2,\quad \alpha=1/10,\quad
 \varepsilon=1/2,\quad \delta=1/100,
 \qquad x_1=\cdots=x_n=4,\] where \(n\) is even and tends to infinity. Grant the proof its correct batch-size calibration, rather than exploiting the different total-\(n\) formula printed in Algorithm 3: \[\begin{aligned}
&s=n/2,\quad c=80\log(4/\delta),\quad
 e^{\varepsilon_L}=\varepsilon^2s/c,\\
&q_n=\frac{1}{1+e^{\varepsilon_L}}
     =\frac{c}{\varepsilon^2s+c}.
\end{aligned}\] Here \(q_n\) is the probability that a zero comparison is reported as one. For all sufficiently large \(n\), \(\varepsilon_L>1\) and the condition \(\varepsilon>16\sqrt{\log(4/\delta)/s}\) in Lemma G.5 holds. The data lie in \([B]\), every user is queried once, and permutation of this constant dataset has no effect.\par
The first midpoint of Algorithm 3 is 2. All its true comparisons are zero. If every one of the \(s\) independent randomized responses is one, the debiased estimate exceeds \(1/2\), so the algorithm sets \(R=1\). Its second midpoint, and returned value, is then 1 regardless of the second batch’s responses. This is an incorrect answer: \(F_X(1)=F_X(2)=0\), so the interval in the paper’s own success predicate misses \((0.4,0.6)\). It is also incorrect under the conventional quantile-value predicate. Shuffling preserves the all-one message multiset. Therefore \[\Pr(\text{failure})\ \geq\ q_n^{n/2}
 =\exp[-\Theta(n\log n)].\] For the confidence sequence \(\beta_L=\exp(-L^2)\), (2.1) with the other parameters fixed prescribes \(n=O(L)\). For every fixed hidden constant and sufficiently large \(L\), the event in (2.5) has probability greater than \(\beta_L\), since \(O(L\log L)=o(L^2)\). Rounding \(n\) to an even integer does not change this conclusion. This is an asymptotic check against every fixed constant, not a numerical comparison to an unknown constant.\par
\#\#\#\# Impact and alternatives.\par
The invalid step links G.3 directly to G.6 and the headline shuffle guarantee. It therefore warrants Major. The construction refutes this protocol’s stated unrestricted confidence dependence, not the existence of every conceivable shuffle protocol with such a rate. Fixed-confidence results, or a suitably restricted joint regime of \(\delta\) and \(\beta\), may survive. Section 2 (p. 4) describes \(\beta=1/\operatorname{poly}(B)\) as the typically interesting regime, but neither Theorem 1.4 nor G.6 imposes that restriction. The counterexample addresses their unrestricted statement; a restricted theorem would need to specify its joint parameter regime. This confidence-dependence counterexample does not refute the intended adaptive LDP upper bound or adaptivity separation.\par
A full correction must retain the tail term and recheck finite-population sampling and parameter regimes. The example proves the necessary condition \(n\log n=\Omega(\log(1/\beta))\) for this calibration. Retaining the Bernstein term gives a conservative local sufficient-batch condition with an additional \(\alpha^{-1}\log(1/\beta')\) term; this review does not certify a complete revised shuffle theorem.\par
\# Implementation and output-contract findings\par
Problem 2 (Minor: the pinned endpoint selector always chooses left).  Algorithm 2, lines 4–5, p. 17, selects the left endpoint when \((q-W(j-1))/w_j\leq q\), otherwise the right endpoint. In `BaySS/bayesian\_learning.py`, the function `get\_interval\_coin\_from\_quantile` passes the inclusive \(W(j)\) to `round\_interval\_to\_coin`. The latter computes \((q-W(j))/w_j-q\).\par
Because the selected interval satisfies \(W(j)\geq q\) and \(w_j>0\), the implemented expression is always nonpositive for \(q>0\). Thus its right-endpoint branch cannot be taken in exact arithmetic.\par
A reachable state demonstrates a substantive query change. Start from three uniform interval weights and take update parameter \(\alpha=1/10\). After the first private bit is zero, the code’s update gives \[(w_1,w_2,w_3)=(4/15,\,1/3,\,2/5).\] At \(q=1/2\), the next selected interval is the second. The intended ratio is \((1/2-4/15)/(1/3)=7/10>1/2\), selecting right; the code’s flag is \(-4/5\), selecting left. The first private outcome has positive probability under randomized response, so this is not an unreachable weight vector.\par
Passing \(W(j)-w_j\) instead supplies the intended left cumulative mass and fixes this branch algebraically. No author code was run. The inspected experimental driver calls this mechanism, but no immutable execution-to-figure link was supplied. A different heuristic could still work empirically; neither a score reduction nor invalid local privacy follows from this query-selection discrepancy. The supported finding is a local implementation/fidelity defect, with empirical magnitude unresolved.\par
Problem 3 (Minor: the stated unary-encoding baseline consumes twice its budget).  Appendix D, pp. 14–15, describes one-hot encoding with independent symmetric randomized response on every bit, each retaining its input with probability \(p=e^\varepsilon/(1+e^\varepsilon)\), and describes the resulting protocol as \(\varepsilon\)-LDP.\par
Two possible values of one user’s data encode as \((1,0)\) and \((0,1)\) at a two-category level. The event that the released vector is \((1,0)\) has probabilities \(p^2\) and \((1-p)^2\), respectively. Their ratio is \[\frac{p^2}{(1-p)^2}=e^{2\varepsilon}>e^\varepsilon
 \qquad(\varepsilon>0).\] Thus the described symmetric mechanism is \(2\varepsilon\)-LDP in the worst case. Per-bit parameter \(\varepsilon/2\) fixes the user-level guarantee exactly and changes this baseline’s rate only by a constant factor.\par
This does not establish a privacy failure in the experiments: the pinned `hierarchical\_mechanism/ldp\_protocol.py` requests `UEClient(...,use\_oue=True)`, an optimized unary-encoding path, rather than the symmetric construction analyzed above. The main proposed algorithm’s single comparison bit is also unaffected. The scope and severity are consequently local.\par
Problem 4 (Major: the formal quantile contract has no legal output on an allowed input).  Section 2, p. 4, requires \([F_X(m),F_X(m+1)]\) to intersect \((q-\alpha,q+\alpha)\), while Definitions 2.2–2.4, p. 5, constrain \(m\) to \([B]=\{1,\ldots,B\}\). The proof of Theorem 1.1, p. 8, returns \(j+1\) after establishing bounds at \(j,j+1\); Appendix C, p. 14, returns \(j\).\par
For \(B=4\), all \(x_i=1\), \(q=1/2\) and \(\alpha=1/10\), every allowed \(m\) has \(F_X(m)=F_X(m+1)=1\). Hence no allowed output meets the stated contract, regardless of sample size or privacy. This is an output-domain defect, not a noise limitation.\par
Two coherent conventions are possible: output a gap index in \(\{0,\ldots,B-1\}\) with \(F_X(0)=0\), or output a quantile value in \([B]\) and use \([F_X(m-1),F_X(m)]\). The latter gives a valid answer \(m=1\) in this example. Definitions 2.2–2.4 supply the success contract for the central empirical, statistical and shuffle guarantees. An allowed dataset for which every legal output fails makes those as-written guarantees unattainable, independently of privacy or sample complexity, and therefore supports Major. A uniform indexing shift can preserve the intended statistical objective and rates, but this possible repair does not validate the original success statement. All algorithms, padding reductions and success evaluators must use the chosen convention consistently; that full conversion and the resulting complete guarantees have not been verified. No impossibility of the intended adaptivity separation is asserted.\par
\# Other checks and limits\par
The first-half random-permutation drift argument in Lemmas 3.1 and 4.1 has the expected martingale structure: the fraction of remaining zeros has conditional mean equal to its current value and increments at most \(1/n\) during the first half of a \(2n\)-element permutation. The stated maximum-deviation bound and union bound support this scoped step. Randomized response maps a probability \(q\) to \((e^\varepsilon q+1-q)/(e^\varepsilon+1)\), with slope \(\tanh(\varepsilon/2)\); the reduction’s small-\(\varepsilon\) scaling is consistent.\par
The main potential argument’s separation into good and bad queries and the information-theoretic lower-bound constructions were inspected, but the sharp external constants and every reduction detail were not independently certified. No issue is asserted merely because those external arguments were not reproduced.\par
Further evidence is needed for the shuffle argument’s use of remaining-user CDFs through all batches, since Lemma 3.1 explicitly covers only the first half. Algorithm 3’s total-\(n\) calibration also differs from the batch-size calibration in the proof. The pinned utility includes a \(\log B\) factor; this does not by itself certify its complete privacy configuration. These pending points add no severity to the demonstrated finding.\par
The experiments disclose use of all users, a tuned update parameter and 200 runs. Their error-bar formula on p. 27 is the standard error of a Bernoulli sample average. No claim of empirical reversal is made. The excluded data and result binaries prevent a provenance-complete reproduction, and author-code execution was prohibited. The arXiv metadata lists only v1, 5 February 2025 [Retained latex: \textbackslash{}cite\{arxiv\}]; no later correction was verified.\par
\# Conclusion\par
Concern | Judgment | Repair / remaining work\par
Shuffle confidence dependence | Major | Complete revised theorem pending\par
Endpoint-selection code | Minor | Local arithmetic repair verified; rerun absent\par
Unary-encoding budget | Minor | Per-bit \(\varepsilon/2\) repair verified\par
Quantile output contract | Major | Consistent contract and dependent guarantees needed\par
Scientific disposition | Pending | No acceptance or author adoption\par
\# Evidence, computation and provenance\par
Original SHA-256: `e81196cf12aba168c18c548a32a4014f2d496db4629d7bd678ba7761b2c60867`. The source manifest, input verification, numbered static source excerpts and their hashes are recorded under `verification/0114/evidence/`. The independent script `independent\_checks.py` uses only standard-library arithmetic and file reading; its output `independent\_checks.json` contains the exact rational endpoint calculation, privacy likelihood ratio and log-domain tail examples. For \(n=20000\), the granted batch calibration gives \(\varepsilon_L=1.65168\ldots\), satisfies G.5 and yields log failure-event probability \(-18270.86898\ldots\); these numerical examples supplement the asymptotic proof.\par
Original pages 17 and 25 were rendered and visually read. Final report pages are separately compiled from the exact final TeX with existing `pdflatex -no-shell-escape` and visually inspected. Build commands, return codes, hashes, visual review and the Lumen save receipt are linked in the accompanying provenance. No author code, installation, author contact or publication occurred.\par
The current assigned AI conversation performed the scientific work, calculations and document verification. `scientific\_decision=pending`; `human\_reviews=[]`. Batch token: `[internal identifier omitted]`.\par
9 A. Aamand et al., Lightweight Protocols for Distributed Private Quantile Estimation, PMLR 267 (2025), 27–58. https://proceedings.mlr.press/v267/aamand25a.html. Assigned published original. Author repository, https://github.com/NynsenFaber/Quantile\_estimation\_with\_adaptive\_LDP, commit `1301f6cdf662befb443dbfa7a94394b6a87bbdcf`. Assigned selective snapshot; static inspection. Same authors, arXiv:2502.02990v1, 5 February 2025. https://arxiv.org/abs/2502.02990. Metadata consulted 24 October 2026.\par
\end{document}
